\section{The promise of one graph}
\label{sec:ch01-promise}

The pitch that got GraphQL into so many companies had nothing to do with
distribution. It was that a client could ask for exactly the data a screen
needed, once, and get back a response shaped like the screen.

\begin{minted}{graphql}
# Sketch: one request for a product page. This is the promise that made
# GraphQL spread, and it is worth keeping in view for the whole chapter.
query ProductPage($sku: String!) {
  product(sku: $sku) {
    title
    description
    price { amount currency }
    availableQuantity
    reviews(first: 3) { nodes { rating body } }
  }
}
\end{minted}

Anyone who has assembled that page from four REST endpoints, each returning a
slightly different notion of a product, each with its own pagination style and
its own opinion about whether \code{null} means missing or zero, understands the
appeal immediately. The mobile team stops asking the backend team for a new
endpoint every sprint. The backend team stops maintaining nine variants of
\code{/product} that differ only in which fields they include. Nobody has to
negotiate a payload.

That is the part everyone remembers. The part that matters more for this book is
the second promise, which is organisational rather than technical. If there is
one schema, there is one place where the company's domain is written down.
\code{Product} means one thing. \code{Money} has one shape. When somebody asks
what data the company has about a customer, the answer is a document that is
guaranteed to be current, because it is the thing serving traffic. Apollo made
this the first of its principles for running GraphQL at scale, arguing that a
company should have \enquote{one unified graph, instead of multiple graphs
created by each team}~\autocite{schmidt2019principled}. The instinct is right,
whatever you make of the source: a single graph is worth more than several
disconnected ones, because the value is in the
joins.\index{one graph principle}

\index{schema!as shared vocabulary}%
For a while, this works beautifully, and it is worth being precise about why,
because the reason is not technical merit. It is that coordination is free when
the graph is small. One team owns the schema. The people who would need to
agree about \code{Product} sit near each other, or at least attend the same
standup. A field gets added in the morning and ships in the afternoon. If
somebody proposes a name that will age badly, the person who will suffer for it
is in the room and says so.

Under those conditions a single GraphQL service is close to the best available
architecture, and I want to be clear that I mean that. Nothing in the following
twenty-seven chapters is an argument that you should have started differently.
The correct first move is one service, one schema, one deployable. The problems
this book solves are the problems of success: more teams, more domains, more
people with a legitimate claim on the same types.

\index{coordination cost}%
What changes is not the code. A schema with four hundred fields is not
meaningfully harder to compile than one with forty. What changes is the number
of people who have to agree before a field can move, and that number grows
faster than the schema does. Ten engineers on one graph is a conversation.
Eighty engineers on one graph is a queue.

Figure~\ref{fig:ch01-one-schema-many-teams} is the shape of that queue. It is
worth looking at for a moment, because it is an unremarkable diagram of an
unremarkable architecture, and everything that goes wrong in the next section
is already visible in it.

\begin{figure}[htbp]
  \centering
  % Many teams, one schema file, one deployable. Referenced from chapter 01.
% Bare tikzpicture: the figure environment, caption and label live at the
% call site. Uses only arrows.meta and positioning (loaded in packages.tex).
\begin{tikzpicture}[
    font=\small,
    team/.style={draw, rounded corners=2pt, minimum width=21mm,
                 minimum height=7mm, align=center},
    core/.style={draw, thick, minimum width=76mm, minimum height=11mm,
                 align=center},
    consumer/.style={draw, dotted, minimum width=21mm, minimum height=7mm,
                     align=center},
    flow/.style={-{Stealth[length=2mm]}, thick},
    note/.style={font=\scriptsize\itshape, align=left}
  ]

  \node[team] (t1) at (-3.45,2.1) {Catalog};
  \node[team] (t2) at (-1.15,2.1) {Orders};
  \node[team] (t3) at ( 1.15,2.1) {Pricing};
  \node[team] (t4) at ( 3.45,2.1) {Accounts};

  \node[core] (schema) at (0,0) {one schema, one deployable};

  \foreach \t in {t1,t2,t3,t4}{%
    \draw[flow] (\t) -- (schema);
  }

  \node[consumer] (web) at (-2.3,-2.1) {web};
  \node[consumer] (ios) at ( 0,-2.1) {iOS};
  \node[consumer] (part) at ( 2.3,-2.1) {partners};

  \foreach \c in {web,ios,part}{%
    \draw[flow] (schema) -- (\c);
  }

  \node[note, anchor=west] at (4.3,1.05)
    {every change is a PR\\against the same file};
  \node[note, anchor=west] at (4.3,-1.05)
    {every change ships\\on the same deploy};

\end{tikzpicture}

  \caption{One schema and one deployable, shared by every team that has
  anything to say about the domain. The arrows into the box are the ones that
  cause trouble: every change is a pull request against the same file, and
  every change ships when the whole service ships.}
  \label{fig:ch01-one-schema-many-teams}
\end{figure}
